\begin{frame}{Soma de prefixos}
    \metroset{block=fill}

    Calculamos um vetor $p$ da seguinte forma:
    \begin{gather}
        p[i] = \begin{cases}
            v[0], \text{ se } i = 0, \\
            v[i] + p[i-1], \text{ caso contrário.}
        \end{cases}
    \end{gather}
    \only<2-4>{
        O que é o valor $p[r] - p[l-1]$. \pause É exatamente
        \begin{align*}
            &v[0] + v[1] + \cdots + v[l-1] + v[l] + \cdots + v[r] - \\
            (&v[0] + v[1] + \cdots + v[l-1]) = \\ \pause
            & v[l] + \cdots + v[r] \ .
        \end{align*}
    }
    
    \onslide<5->{
        \begin{block}{Solução linear}
            \footnotesize
            \inputminted{cpp}{codigos/psum.cpp}    
        \end{block}
        Complexidade: $\pause O(n + q)$.
    }
\end{frame}

